\chapter*{Résumé}
\addcontentsline{toc}{chapter}{Résumé}
\markboth{Résumé}{Résumé}

La plateforme chimique de Jorf Lasfar produit de l'acide phosphorique sur six lignes
d'attaque-filtration, le concentre sur cinq lignes, puis le purifie par trois procédés
avant de le distribuer à neuf consommateurs. Chaque jour, un ingénieur d'exploitation doit
décider quelles quantités décadmier, quelles quantités clarifier, quels transferts effectuer
entre ateliers et quelle ligne sert quel client. Ces décisions sont couplées, contraintes
par un réseau de tuyauterie incomplet, et prises aujourd'hui au jugement.

Ce travail formalise ce problème en un \textbf{programme linéaire mixte en nombres entiers}
mono-période et en fournit une implémentation logicielle complète, testée et validée.

Trois apports méthodologiques structurent le rapport. Premièrement, l'identification de
l'unité de compte réelle du système --- la \textbf{tonne de \PdeuxOcinq{}} et non la tonne
d'acide marchand --- information absente de toute la documentation fournie, établie ici par
trois preuves indépendantes, et sans laquelle l'intégralité des bilans matière serait
erronée. Deuxièmement, la démonstration que la fonction objectif spécifiée dans le cahier
des charges est \textbf{mathématiquement dégénérée} : maximiser le total livré sous
contrainte de satisfaction exacte de la demande produit un objectif constant sur tout le
domaine réalisable. Un objectif \textbf{lexicographique à trois niveaux} est proposé, dont
le premier niveau est strictement équivalent à l'objectif demandé et le généralise au cas
infaisable. Troisièmement, un audit systématique des spécifications a mis en évidence
\textbf{onze anomalies}, dont quatre erreurs arithmétiques atteignant \SI{59}{\percent}
d'écart ; chacune est prouvée et arbitrée selon une hiérarchie de sources établie a priori.

Sur le scénario de référence, le modèle atteint le statut optimal en moins de \SI{0.1}{\second} et
satisfait \SI{100}{\percent} de la demande. Un validateur écrit indépendamment du modèle,
à partir des seules équations physiques, confirme la solution par 86 contrôles.

Le résultat le plus utile pour l'exploitation n'est cependant pas le plan lui-même, mais le
diagnostic qui l'accompagne : le modèle révèle deux violations structurelles --- le bac
central IR11 déborde de \SI{1016}{\tonne} par jour et le stock d'une ligne reste sous son
seuil de sécurité --- qu'\emph{aucune} décision d'exploitation ne peut corriger. Une analyse
de sensibilité en établit la cause unique et le remède chiffré : réduire le nombre
d'échelons affectés à la cocristallisation résorbe intégralement les deux violations, sans
investissement et sans dégrader le service rendu. Un mode optionnel, dans lequel cette
affectation devient une variable de décision, confirme ce diagnostic --- sous la réserve,
explicitement portée, qu'un modèle mono-période n'a aucune raison de reconstituer un stock
central.

\vspace{0.8cm}
\noindent\textbf{Mots-clés} : recherche opérationnelle, programmation linéaire en nombres
entiers, optimisation lexicographique, bilan matière, acide phosphorique, planification de
production, OCP.

\cleardoublepage
